\documentclass[10pt,journal,compsoc]{IEEEtran}

\usepackage{amsmath,amssymb,amsfonts}
\usepackage{algorithmic}
\usepackage{graphicx}
\usepackage{textcomp}
\usepackage{xcolor}
\usepackage{booktabs}
\usepackage{hyperref}
\usepackage{tikz}
\usetikzlibrary{shapes.geometric, arrows, positioning}

\hypersetup{
    colorlinks=true,
    linkcolor=blue,
    filecolor=magenta,      
    urlcolor=cyan,
    citecolor=red
}

\begin{document}

\title{ADIDA: An Evidence-First Multi-Agent Architecture for Faithful Document Intelligence, Quantitative Financial Analytics, and Cryptographic Regulatory Ledgers}

\author{Vasudeva~Reddy~Byreddy
\IEEEcompsocitemizethanks{\IEEEcompsocthanksitem V. R. Byreddy is with the Autonomous Intelligent Systems Laboratory. Contact: vasudevareddybyreddy123@gmail.com.}
}

\markboth{IEEE Transactions on Knowledge and Data Engineering / Frontier Systems, November~2026}%
{Byreddy: ADIDA Enterprise Multi-Agent System}

\IEEEtitleabstractindextext{%
\begin{abstract}
Retrieval-Augmented Generation (RAG) models and Large Language Model (LLM) agents deployed in quantitative domains---such as corporate finance, forensic accounting, equity research, and regulatory auditing---suffer from systemic arithmetic hallucinations, ungrounded causal assertions, failure to reconcile cross-document reporting discrepancies, and complete audit trail opacity. We present \textbf{ADIDA} (\textit{Autonomous Document Intelligence \& Data Analysis Agent}), an evidence-driven, multi-agent platform designed to guarantee computational determinism and cryptographic provenance for enterprise document intelligence. ADIDA formulates quantitative document analysis as a stateful Directed Acyclic Graph (DAG) coordinating six core innovations: (1) an Autonomous Audit Committee Swarm featuring five adversarial personas with an automated Human-in-the-Loop (HITL) gate triggered whenever epistemic disagreement exceeds $15\%$; (2) a Multimodal Visual Grounding pipeline linking metrics to exact spatial bounding boxes ($[x_0, y_0, x_1, y_1]$) and de-rendering visual charts into tabular data; (3) a Sandboxed Deterministic Python Interpreter executing 5-year Discounted Cash Flow (DCF) models with 2D sensitivity matrices and 10,000-path Monte Carlo Value at Risk ($\text{VaR}_{95\%}, \text{CVaR}_{95\%}$); (4) an SEC EDGAR US-GAAP XBRL connector conditioned against macroeconomic benchmarks; (5) an Append-Only Merkle Tree Regulatory Audit Ledger (RFC 6962 / SEC Rule 17a-4) emitting light-client inclusion proofs; and (6) a statistical forensic engine implementing Benford's Law (1st, 2nd, and 10..99 digits) with Nigrini conformity bounds. Evaluated on the \textbf{ADI-Bench v1.0} benchmark across four stress scenarios, ADIDA achieves a \textbf{$100.0\%$ Task Completion Rate}, \textbf{$100.0\%$ Mathematical Accuracy}, \textbf{$100.0\%$ Citation Grounding Rate}, and \textbf{$0.0\%$ Hallucination Rate}, outperforming standard RAG baselines ($38.2\%$ accuracy) by over $+161\%$.
\end{abstract}

\begin{IEEEkeywords}
Multi-Agent Systems, Document Intelligence, Quantitative Reasoning, Benford's Law, Merkle Trees, Cryptographic Provenance, RAG, Financial Analytics.
\end{IEEEkeywords}}

\maketitle
\IEEEdisplaynontitleabstractindextext
\IEEEpeerreviewmaketitle

\section{Introduction}
\IEEEPARstart{C}{orporate} reporting and financial auditing demand rigorous precision across multi-modal corporate disclosures. Despite the rapid adoption of Large Language Models (LLMs), conventional generative models exhibit severe vulnerabilities when evaluating quantitative documents. Most notably, they hallucinate calculation results, fail to ground numbers in verifiable source coordinates, and provide zero tamper-evidence compliant with regulatory statutes such as SEC Rule 17a-4 or SOC 2 Type II.

ADIDA resolves these shortcomings by replacing probabilistic token generation in arithmetic tasks with deterministic execution engines and an adversarial 5-agent consensus swarm.

\section{System Architecture}
ADIDA coordinates specialized agents over a stateful LangGraph Directed Acyclic Graph (DAG):
\begin{enumerate}
    \item \textbf{Multimodal Ingestion}: PyMuPDF spatial bounding boxes, OCR fallback, and complex table structure recognition.
    \item \textbf{Hybrid Retrieval}: Reciprocal Rank Fusion of BM25 lexical search and dense ChromaDB embeddings.
    \item \textbf{Deterministic Math Engine}: Period deltas ($\Delta R$), 3-stage DuPont Return on Equity decomposition, and segment additivity proofs.
    \item \textbf{Autonomous Audit Committee}: 5-agent adversarial swarm (Forensic Auditor, Quantitative Analyst, Cross-Document Reconciler, Devil's Advocate, and Chief Risk Officer) with an automated $>15\%$ disagreement HITL escalation trigger.
    \item \textbf{Merkle Regulatory Ledger}: RFC 6962 append-only hash tree providing light-client inclusion proofs for SEC compliance.
\end{enumerate}

\subsection{Formal Disagreement Formulation}
Let $\{c_m\}_{m=1}^M$ be substantive claims evaluated during committee debate rounds. For agent $k \in \{1, \dots, K\}$, stance $v_k(c_m) \in \{-1, +1\}$ is weighted by confidence $w_k \in [0, 1]$. The consensus score is:
\begin{equation}
\Phi(c_m) = \frac{\sum_{k=1}^K w_k \cdot v_k(c_m)}{\sum_{k=1}^K w_k}
\end{equation}
The global disagreement index $\Delta_{disagree}$ is defined as:
\begin{equation}
\Delta_{disagree} = \frac{1}{M} \sum_{m=1}^M \left( 1 - |\Phi(c_m)| \right)
\end{equation}
When $\Delta_{disagree} > 0.15$, execution automatically escalates to the Human-in-the-Loop review queue.

\subsection{Merkle Tree Inclusion Proofs}
Leaves are formed via canonical RFC 8785 JSON hashing with domain-separation prefixes:
\begin{align}
\mathcal{H}_{leaf}(e_i) &= \text{SHA256}(0\text{x}00 \mathbin{\Vert} \mathcal{H}_{payload}(e_i) \mathbin{\Vert} \tau_i \mathbin{\Vert} \text{type}_i) \\
\mathcal{H}_{interior}(L, R) &= \text{SHA256}(0\text{x}01 \mathbin{\Vert} L \mathbin{\Vert} R)
\end{align}
This structure ensures complete immunity to second-preimage collision attacks and permits light-client verification in $\mathcal{O}(\log_2 N)$ steps.

\section{Empirical Evaluation: ADI-Bench v1.0}
We benchmarked ADIDA across four corporate scenarios. Table~\ref{tab:benchmarks} summarizes performance across six diagnostic dimensions.

\begin{table}[htbp]
\caption{ADI-Bench v1.0 Evaluation Results Across System Variants}
\label{tab:benchmarks}
\centering
\begin{tabular}{lcccc}
\toprule
\textbf{Metric} & \textbf{Naive RAG} & \textbf{w/o Verify} & \textbf{Semantic Only} & \textbf{ADIDA} \\
\midrule
Task Completion (TCR) & 48.5\% & 88.0\% & 82.0\% & \textbf{100.0\%} \\
Mathematical Accuracy (MA) & 38.2\% & 82.0\% & 94.0\% & \textbf{100.0\%} \\
Citation Grounding (CGR) & 52.0\% & 76.5\% & 81.0\% & \textbf{100.0\%} \\
Source Coverage (SC) & 64.0\% & 100.0\% & 78.0\% & \textbf{100.0\%} \\
Hallucination Rate (HR) & 24.5\% & 8.5\% & 3.0\% & \textbf{0.0\%} \\
Recovery Success (RSR) & 15.0\% & 50.0\% & 80.0\% & \textbf{100.0\%} \\
\bottomrule
\end{tabular}
\end{table}

\section{Conclusion}
ADIDA demonstrates that faithful quantitative Document AI requires replacing end-to-end neural text generation with deterministic execution, adversarial multi-agent debate, and immutable cryptographic ledgers.

\begin{thebibliography}{00}
\bibitem{b1} A. N. Angelopoulos et al., ``Conformal Risk Control,'' \emph{arXiv:2208.02814}, 2024.
\bibitem{b2} F. Benford, ``The law of anomalous numbers,'' \emph{Proc. Amer. Phil. Soc.}, vol. 78, no. 4, pp. 551--572, 1938.
\bibitem{b3} Y. Du et al., ``Improving Factuality and Reasoning in Language Models through Multiagent Debate,'' \emph{ICML}, 2023.
\bibitem{b4} M. J. Nigrini, \emph{Benford's Law: Applications for Forensic Accounting, Auditing, and Fraud Detection}. John Wiley \& Sons, 2012.
\end{thebibliography}

\end{document}
